\section{Architekturentscheidungen}

Die folgenden Entscheidungen prägen den Aufbau und die Implementierung
des Systems. Neben der Begründung werden auch die daraus entstehenden
Nachteile und zusätzlichen Aufwände betrachtet.

    {
    \renewcommand{\arraystretch}{1.4}

    \begin{tabularx}{\linewidth}{@{}
            >{\bfseries\raggedright\arraybackslash}p{0.09\linewidth}
            >{\raggedright\arraybackslash}p{0.22\linewidth}
            >{\raggedright\arraybackslash}X
            >{\raggedright\arraybackslash}X
        @{}}

\toprule
\rowcolor{tableHeaderColor}
Nr. & Entscheidung & Begründung & Konsequenz \\
\midrule

AE 1 &
Fachliche Aufteilung in eigenständige Services &
Authentifizierung, Benutzerverwaltung, Produkte, Bestellungen und
Benachrichtigungen besitzen getrennte Verantwortlichkeiten und sollen
unabhängig entwickelt werden können. &
Die Services können getrennt entwickelt und betrieben werden.
Dafür entstehen zusätzlicher Deployment-Aufwand, Netzwerkkommunikation
und verteilte Fehlerfälle. \\

\rowcolor{tableRowColor}
AE 2 &
Zentrales API Gateway &
Routing und JWT-Prüfung sollen nicht in jedem Service einzeln
implementiert werden. &
Das Gateway bildet einen einheitlichen Zugang, verursacht aber einen
zusätzlichen Netzwerkhop und ist ein zentraler Ausfallpunkt. \\

AE 3 &
Eigene Datenbank je Service &
Jeder Service soll alleiniger Eigentümer seiner Daten und seines
Datenmodells sein. &
Die Services bleiben hinsichtlich ihrer Datenhaltung unabhängig.
Serviceübergreifende Fremdschlüssel und gemeinsame Transaktionen sind
jedoch nicht möglich. \\

AE 4 &
Kompensation statt verteilter Transaktion &
Product Service und Order Service besitzen getrennte Datenbanken,
sodass keine gemeinsame Datenbanktransaktion verwendet werden kann. &
Bereits geänderte Produkte werden bei einem Fehler durch Gegenaufrufe
wieder freigegeben. Diese Rückabwicklung erfordert zusätzliche
Programmlogik. \\

\rowcolor{tableRowColor}
AE 5 &
Zustandslose Product-Service-Replikate &
Produktanfragen sollen exemplarisch auf mehrere Instanzen verteilt
werden können. &
Drei Product-Service-Instanzen können über Nginx angesprochen werden.
Die gemeinsam verwendete Produktdatenbank bleibt jedoch ein möglicher
Engpass. \\

\bottomrule
    \end{tabularx}
}